\section{Missing Proofs of Section~\ref{sec:k-seq}}\label{sec:k-seq-appendix}

In this section, we provide the missing proofs of lemmas and claims made in Section~\ref{sec:k-seq}.

\begin{lemma}\label{lemma:augment}
    For any partitioning $\mcC$ of the root square $\cell^*$ of $\mcH$, any feasible extended matching $M^\mcC=(M, B^\mcC), y(\cdot)$ on $\mcG_\sigma$, and any admissible alternating or augmenting path $P$, the matching obtained after updating $M^\mcC$ along $P$ remains feasible.
\end{lemma}
\begin{proof}
    For any edge $(a,b)\in P$, due to the admissibility of the edge, $y(b)-y(a)=\distance{a}{b}$. If $(a,b)\notin M$ prior to augmentation, it is a matching edge after the augmentation and Condition~\eqref{eq:dualfeasibility-matching-time} holds for $(a,b)$. Otherwise, $(a,b)\in M$ prior to augmentation, it is a non-matching edge after the augmentation, and Condition~\eqref{eq:dualfeasibility-non-matching-time} holds for $(a,b)$. Note that all dual weights remain unchanged and consequently, Conditions~\eqref{eq:dualfeasibility-b}--\eqref{eq:dualfeasibility-a-free} remain satisfied.
\end{proof}


% \myparagraph{Proof of Invariants (I1) and (I2).} We begin by showing that (I1) and (I2) continue to be satisfied after locally processing $\cell$. We then show the same for the global search. 

\subsection{Correctness of the Extended Hungarian Search Procedure}

\hungarianProp*
% \begin{restatable}{lemma}{globalProp}\label{lemma:global}
%     For any cell $\cell\in\mcC$, let $M_\cell, y(\cdot)$ denote a $\cell$-optimal matching on $\mcG_\cell$. Then, assuming invariants (I1) and (I2), after executing the extended Hungarian search procedure on $\cell$, the matching remains $\cell$-optimal. Furthermore, the dual weight of all active free points inside $\cell$ will be $\ymax$.
% \end{restatable} 



Let $y(\cdot)$ (resp. $y'(\cdot)$) denote the dual weights of the points after (resp. before) the update duals step of the extended Hungarian search procedure. To prove this lemma, we first show that after the update duals step, the extended matching $M^\mcC, y(\cdot)$ is feasible, $y(v)\le \ymax$ for all $v\in A_\cell\cup B_\cell$, and $y(b_f)=\ymax$ for all free points $b_f\in B_\cell$. We then show that the augmenting path $P$ computed by the procedure is admissible and conclude that $P$ is a minimum net-cost augmenting path inside $\cell$. We then use Lemma~\ref{lemma:augment} to show that after augmentation, the extended matching $M^\mcC, y(\cdot)$ remains feasible. Finally, we show that the updated key of $\cell$ is the smallest net-cost of augmenting paths inside $\cell$ and is at least $\ymax$.

    \myparagraph{Feasibility of points.}
    \begin{enumerate}
        \item For any point $b\in B$:
        \begin{itemize}
            \item if $b\in B\setminus B_\cell$ is outside of $\cell$, then $y(b)=y'(b)$; therefore, Conditions~\eqref{eq:dualfeasibility-b} and~\eqref{eq:dualfeasibility-b-admissible} remains satisfied.
            \item Otherwise, $b\in B_\cell$ and by definition, $\kappa\le \kappa_b+s(b)$. Therefore, Condition~\eqref{eq:dualfeasibility-b} holds for $b$ since 
            \begin{equation}\label{eq:global-hung-2}
                y(b) = y'(b) + \kappa - \kappa_b \le y'(b) + s(b) =\distancetocell{b}{\cell}. 
            \end{equation}
        \end{itemize}
        \item For any free point $a\in A_F$:
        \begin{itemize}
            \item if $a\in A\setminus A_\cell$ is a free point outside of $\cell$, then $y(a)=y'(a)=0$ and Condition~\eqref{eq:dualfeasibility-a-free} remains true. 
            \item Otherwise, $a\in A_\cell$ and by definition, $\kappa_a\ge \kappa$; hence, the procedure does not update the dual weight of $a$, i.e., the dual weight of $a$ remains $0$, satisfying Condition~\eqref{eq:dualfeasibility-a-free}. 
        \end{itemize}
        
    \end{enumerate}

    \myparagraph{Feasibility of edges.} For any edge $(a,b)\in E$, let $s(a,b)$ denote the slack of $(a,b)$ before updating the dual weights. 
    \begin{enumerate}
        \item if $a\in A\setminus A_\cell$ and $b\in B\setminus B_\cell$, then $y(a)=y'(a)$ and $y(b)=y'(b)$ and Conditions~\eqref{eq:dualfeasibility-non-matching-time} and~\eqref{eq:dualfeasibility-matching-time} remains satisfied for $(a,b)$.
        
        \item if $a\in A_\cell$ and $b\in B\setminus B_\cell$, then $y(b)=y'(b)$ and $y(a)\ge y'(a)$, since the extended Hungarian search procedure does not decrease the dual weight of any point inside $\cell$; hence, $y(b)-y(a)\le y'(b)-y'(a)\le \distance{a}{b}$ and Condition~\eqref{eq:dualfeasibility-non-matching-time} is satisfied. (Note that by Lemma~\ref{lemma:no-cross}, the edge $(a,b)$ is a non-matching edge).

        \item if $a\in A\setminus A_\cell$ and $b\in B_\cell$, then $y(b)\le \distance{b}{\mcC}\le \distance{a}{b}$; hence, $y(b)-y(a)\le y(b)\le \distance{a}{b}$ and Condition~\eqref{eq:dualfeasibility-non-matching-time} is satisfied. (Note that by Lemma~\ref{lemma:no-cross}, the edge $(a,b)$ is a non-matching edge).

        \item if $a\in A_\cell$ and $b\in B_\cell$: 
        \begin{itemize}
            \item If $(a,b)\in M$ is a matching edge, then $\kappa_b = \kappa_a$ since the only edge directed to $b$ in the residual graph is the zero-slack edge $(a,b)$. Thus, Condition~\eqref{eq:dualfeasibility-matching-time} holds since
            \begin{equation*}
            y(b) - y(a) = (y'(b)+\kappa-\kappa_b) - (y'(a)+\kappa-\kappa_a) = y'(b)-y'(a) = \distance{a}{b}.
            \end{equation*}
            \item Otherwise, $(a,b)$ is a non-matching edge and $\kappa_a\le \kappa_b + s(b,a)$; therefore, Condition~\eqref{eq:dualfeasibility-non-matching-time} holds since
            \begin{equation}\label{eq:global-hung-1}
                y(b)-y(a) = (y'(b)+\kappa-\kappa_b) - (y'(a)+\kappa-\kappa_a) \le y'(b)-y'(a) + s(b,a) = \distance{a}{b}.
            \end{equation}
        \end{itemize}
    \end{enumerate}
    
    


\myparagraph{Maximum dual weight.} Let $y_{\cell}:=\max_{b\in B_\cell}y'(b)$. By invariant (I2) prior to the extended Hungarian search procedure, for all free points $b\in B_\cell$, $y'(b)=y_{\cell}$. By the construction of the residual graph, $\kappa_b=y'(b)=y_\cell$ for all free points $b\in B_\cell$ and $\kappa_b\ge y_\cell$ for all points $b\in B_\cell$. Therefore, for any free point $b\in B_\cell$, 
\[y(b) = y'(b)+\kappa-\kappa_{b} = \kappa=\ymax,\]
where the last equality holds since $\kappa$ is the net-cost of the minimum net-cost path inside $\cell$, which is $\ymax$.
Furthermore, for any point $b'\in B_\cell$,
\[y(b) = y'(b)+\kappa-\kappa_{b} \le y'(b)+\kappa-y_\cell \le \kappa=\ymax.\]
Note that for any free point $a\in A_\cell$, $y(a)=0$ and for any matched point $a\in A_\cell$, if $a$ is matched to a point $b\in B_\cell$, then by Condition~\eqref{eq:dualfeasibility-matching-time}, $y(a)=y(b)-\distance{a}{b}\le y(b)\le \ymax$. 
Hence, we conclude $y(v)\le \ymax$ for all points $v\in A_\cell\cup B_\cell$ and $y(b_f)=\ymax$ for all free points $b_f\in B_\cell$ after the dual updates. 



\myparagraph{Net-cost of $P$.} To prove that $P$ is a minimum net-cost augmenting path, we show that $P$ is an admissible augmenting path. Consequently, if $b\in B_\cell$ is the free endpoint of $P$, by Corollary~\ref{cor:slackcost}, $\phi(P)=y(b)=\ymax$. Note that for all other augmenting paths $P'$ from a free point $b'\in B_\cell$, from Lemmas~\ref{lem:slackcost}, $\phi(P')\ge y(b')=\ymax$, and therefore, $P$ would be a minimum net-cost augmenting path inside $\cell$.

For each edge $(u,v)\in P$, $\kappa_v = \kappa_u + s(u,v)$, since $(u,v)$ is an edge of the shortest path tree of the residual graph. Plugging into Equation~\eqref{eq:global-hung-1}, for each non-matching edge $(b,a)$, $y(b)-y(a)=\distance{a}{b}$ and the edge $(b,a)$ is admissible. Furthermore, if $\kappa=\kappa_b+s(b)$ for a point $b\in B$, then by Equation~\eqref{eq:global-hung-2}, $y(b)=\distancetocell{b}{\cell}$ and the point $b$ would be a zero-slack point. Therefore, the path $P$ computed by the algorithm would be an admissible path from a free point $b\in B_\cell$ to either a free point $a\in A_\cell$ or a zero-slack point $b'\in B_\cell$, i.e., $P$ is an admissible augmenting path. 

From Lemma~\ref{lemma:augment}, the extended matching $M^\mcC, y(\cdot)$ obtained after augmenting the matching $M^\mcC$ along $P$ remains feasible.

\myparagraph{Updated key.} Note that the extended matching $M^\mcC, y(\cdot)$ after the augmentation step is feasible, $y(v)\le \ymax$ for all vertices $v\in A_\cell\cup B_\cell$, and $y(b_f)=\ymax$ for all free points $b_f\in B_\cell$. Let $\kappa_v$, for each $v\in A_\cell\cup B_\cell$, denote the distances computed in the update key step of the extended Hungarian search procedure. By Lemma~\ref{lem:slackcost}, for any augmenting path $P$ from a free point $b\in B_\cell$ to a free point $a\in A_\cell$, the net-cost of $P$ is 
\[\phi(P)=y(b) +\sum_{(a'',b'')\in P}s(a'',b'') \ge \kappa_a,\] where the last inequality holds by the construction of the residual graph and the definition of $\kappa_a$. In this case, the equality happens when $P$ is a shortest path from $s$ to $a$. Similarly, for any augmenting path $P$ from a free point $b\in B_\cell$ to a point $b'\in B_\cell$, the net-cost of $P$ is 
\[\phi(P)=y(b) + s(b') +\sum_{(a'',b'')\in P}s(a'',b'') \ge \kappa_{b'} + s(b'),\]
and we get an equality if $P$ is a shortest path from the source vertex $s$ to $b'$.
Therefore, the updated key of $\cell$, i.e., $\kappa_\cell=\min\{\min_{a\in A_\cell^F}\kappa_a, \min_{b\in B_\cell}\kappa_b+s(b)\}$ correctly computes the net-cost of the minimum net-cost augmenting path inside $\cell$. Finally, note that for any augmenting path $P$ from a free point $b\in B_\cell$, by Lemma~\ref{lem:slackcost}, $\phi(P)\ge y(b)=\ymax$ and the key of $\cell$ would be at least $\ymax$.



\subsection{Correctness of the Merge Procedure}\label{subsec:merge_correctness}

\mergeProp*
Let $\hat{M}^\mcC=(M, \hat{B}^\mcC), \hat{y}(\cdot)$ denote the feasible extended matching maintained by our algorithm before the execution of the merge procedure. Let $B^\mcC_\cell$ denote the subset of points in $\hat{B}^\mcC$ that are matched to the divider $\Gamma_\cell$ of $\cell$, and let $M^\mcC=(M, B^\mcC = \hat{B}^\mcC\setminus B^\mcC_\cell)$ be the extended matching after adding the boundary-matched points in $B_\cell^\mcC$ to the free points. By Lemma~\ref{lemma:combination}, the matching $M^\mcC, \hat{y}(\cdot)$ is feasible. Next, we show that, given a feasible extended matching $M^\mcC, y'(\cdot)$, after one iteration of the while-loop in the merge step, the matching remains feasible and the dual weights of points in $\cell$ are at most $\ymax$.

Let $y(\cdot)$ (resp. $y'(\cdot)$) denote the dual weights of the points after (resp. before) the execution of one iteration of the merge step. We first show that $M^\mcC, y(\cdot)$ is a feasible extended matching where $y(v)\le \ymax$ for all $v\in A_\cell\cup B_\cell$. We then show that the path $P$ is admissible and use Lemma~\ref{lemma:augment} to show that after augmentation, the extended matching $M^\mcC, y(\cdot)$ remains feasible.

\myparagraph{Feasibility of points.}
    \begin{enumerate}
        \item For any point $b\in B$:
        \begin{itemize}
            \item if $b\in B\setminus B_\cell$ is outside of $\cell$, then $y(b)=y'(b)$; therefore, Conditions~\eqref{eq:dualfeasibility-b} and~\eqref{eq:dualfeasibility-b-admissible} remains satisfied.
            \item Otherwise, $b\in B_\cell$ and $\kappa\le \kappa_b+s(b)$. Therefore, Condition~\eqref{eq:dualfeasibility-b} holds for $b$ since 
            \begin{equation}\label{eq:local-hung-2}
                y(b) = y'(b) + \kappa - \kappa_b \le y'(b) + s(b) =\distancetocell{b}{\cell}. 
            \end{equation}
        \end{itemize}
        \item For any free point $a\in A_F$:
        \begin{itemize}
            \item if $a\in A\setminus A_\cell$ is a free point outside of $\cell$, then $y(a)=y'(a)=0$ and Condition~\eqref{eq:dualfeasibility-a-free} remains true. 
            \item Otherwise, $a\in A_\cell$ and $\kappa_a\ge \kappa$; therefore, the procedure does not update the dual weight of $a$, i.e., the dual weight of $a$ remains $0$, satisfying Condition~\eqref{eq:dualfeasibility-a-free}. 
        \end{itemize}
        
    \end{enumerate}

    \myparagraph{Feasibility of edges.} For any edge $(a,b)\in E$, let $s(a,b)$ denote the slack of $(a,b)$ with respect to $y'(\cdot)$. 
    \begin{enumerate}
        \item if $a\in A\setminus A_\cell$ and $b\in B\setminus B_\cell$, then $y(a)=y'(a)$ and $y(b)=y'(b)$ and therefore, Conditions~\eqref{eq:dualfeasibility-non-matching-time} and~\eqref{eq:dualfeasibility-matching-time} remains satisfied for $(a,b)$.
        
        \item if $a\in A_\cell$ and $b\in B\setminus B_\cell$, then $y(b)=y'(b)$ and $y(a)\ge y'(a)$, since the merge procedure does not decrease the dual weight of any point inside $\cell$; hence, $y(b)-y(a)\le y'(b)-y'(a)\le \distance{a}{b}$ and Condition~\eqref{eq:dualfeasibility-non-matching-time} is satisfied (Note that by Lemma~\ref{lemma:no-cross}, the edge $(a,b)$ is a non-matching edge).

        \item if $a\in A\setminus A_\cell$ and $b\in B_\cell$, then $y(b)\le \distance{b}{\mcC}\le \distance{a}{b}$; hence, $y(b)-y(a)\le y(b)\le \distance{a}{b}$ and Conditions~\eqref{eq:dualfeasibility-non-matching-time} is satisfied (Note that by Lemma~\ref{lemma:no-cross}, the edge $(a,b)$ is a non-matching edge).

        \item if $a\in  A_\cell$ and $b\in B_\cell$: 
        \begin{itemize}
            \item If $(a,b)\in M$ is a matching edge, then $\kappa_b = \kappa_a$ since the only edge directed to $b$ in the residual graph is the zero-slack edge $(a,b)$. Thus, Condition~\eqref{eq:dualfeasibility-matching-time} holds since
            \begin{equation*}
            y(b) - y(a) = (y'(b)+\kappa-\kappa_b) - (y'(a)+\kappa-\kappa_a) = y'(b)-y'(a) = \distance{a}{b}.
            \end{equation*}
            \item Otherwise, $(a,b)$ is a non-matching edge and since there is a directed edge from $b$ to $a$, $\kappa_a\le \kappa_b + s(b,a)$, and Condition~\eqref{eq:dualfeasibility-non-matching-time} holds since
            \begin{equation}\label{eq:local-hung-1}
                y(b)-y(a) = (y'(b)+\kappa-\kappa_b) - (y'(a)+\kappa-\kappa_a) \le y'(b)-y'(a) + s(b,a) = \distance{a}{b}.
            \end{equation}
        \end{itemize}
    \end{enumerate}

    \myparagraph{Maximum dual weight.} Next, we show that $y(v)\le \ymax$ for all $v\in A_\cell\cup B_\cell$. Note that for any point $b\in B_\cell$, by the definition of $\kappa$, we have $\kappa\le \kappa_b+\ymax-y'(b)$. Therefore, if $\kappa_b<\kappa$, then $y(b) = y'(b)-\kappa_b+\kappa\le \ymax$. Otherwise, $\kappa_b\ge \kappa$ and $y(b)=y'(b)\le \ymax$. Furthermore, for all free points $a\in A_\cell$, by Condition~\eqref{eq:dualfeasibility-a-free}, $y(a)=0$ and for all matched points $a\in A_\cell$, if $a$ is matched to a point $b\in B_\cell$, then $y(a)=y(b)-\distance{a}{b}\le y(b)\le \ymax$.
    


    \myparagraph{Net-cost of $P$.} Next, we show that the path $P$ computed in an iteration of the merge step is either (i) an admissible augmenting path, or (ii) an admissible alternating path from a free point $b\in B_\cell$ to a point $b'\in B_\cell$ with $y(b')=\ymax$. Then, from Lemma~\ref{lemma:augment}, the extended matching obtained after updating $M^\mcC, y(\cdot)$ along $P$ is feasible.
    
    For each non-matching edge $(b,a)\in P$, since $(b,a)$ is in the Dijkstra's shortest path tree, $\kappa_a = \kappa_b + s(b,a)$. Plugging into Equation~\eqref{eq:local-hung-1}, for each non-matching edge $(b,a)$, $y(b)-y(a)=\distance{a}{b}$, and the edge $(b,a)$ is admissible (i.e., all edges of $P$ are admissible).
    \begin{itemize}
        \item If $P$ is a path that ends at a free point $a\in A_\cell$ ($\kappa$ is determined by the first term in the RHS of Equation~\eqref{eq:local-kappa}), then $P$ is an admissible augmenting path, 
        \item otherwise, if $P$ is a path that ends at a point $b'\in B_\cell$ with $\kappa=\kappa_{b'}+s(b')$ ($\kappa$ is determined by the second term in the RHS of Equation~\eqref{eq:local-kappa}), then from Equation~\eqref{eq:local-hung-2}, $y(b') = y'(b')+s(b') = \distance{b'}{\mcC}$, and $P$ is an admissible augmenting path, and
        \item otherwise, $P$ is a path that ends at a point $b'\in B_\cell$ with $\kappa=\kappa_{b'}+\ymax-y'(b')$, and $P$ is an admissible alternating path.
    \end{itemize}
    Therefore, $P$ is an admissible alternating or augmenting path and the matching obtained by updating $M^\mcC$ along $P$ remains feasible (Lemma~\ref{lemma:augment}). 
    
    \myparagraph{Termination.} Let $P$ be a path from a free point $b\in B_\cell$ to a point $u\in \freeofcell{A}{\cell}\cup B_\cell$.
    \begin{itemize}
        \item If $u\in\freeofcell{A}{\cell}$, then $P$ is an admissible augmenting path and $P$ is in case (i).
        \item Otherwise, if $u\in B_\cell$ and $\kappa=\kappa_u+s(u)$, then $P$ is an admissible augmenting path and $P$ is in case (ii).
        \item Otherwise, $u\in B_\cell$ and $\kappa=\kappa_u+\ymax - y'(u)$. In this case, $P$ is an admissible alternating path from $b$ to a free point $b'\in B_\cell$ with $y'(b') < \ymax$ and $y(b')=\ymax$. 
    \end{itemize}
    In either case, the number of free points $b\in B_\cell$ with $y(b)<\ymax$ reduces by one, and the merge step terminates.

    \myparagraph{Dual weight of free points.} Note that the while-loop terminates when there are no free points $b_f\in B_\cell$ with $y(b_f)<\ymax$, whereas, as discussed above, all points will have a dual weight at most $\ymax$, i.e., the dual weight of each free point $b_f\in B_\cell$ after the termination of the while-loop is $\ymax$.

    
\myparagraph{Updated key.} Finally, we show that the updated key of $\cell$ denotes the net-cost of the minimum net-cost augmenting path inside $\cell$. Note that the extended matching $M^\mcC, y(\cdot)$ after the termination of the while-loop is feasible, $y(v)\le \ymax$ for all vertices $v\in A_\cell\cup B_\cell$, and $y(b_f)=\ymax$ for all free points $b_f\in B_\cell$. Let $\kappa_v$, for each $v\in A_\cell\cup B_\cell$, denote the distances computed in the update key step of the extended Hungarian search procedure. By Lemma~\ref{lem:slackcost}, for any augmenting path $P$ from a free point $b\in B_\cell$ to a free point $a\in A_\cell$, the net-cost of $P$ is 
\[\phi(P)=y(b) +\sum_{(a'',b'')\in P}s(a'',b'') \ge \kappa_a,\] where the last inequality holds by the construction of the residual graph and we get an equality of $P$ is the shortest path from $s$ to $a$. Similarly, for any augmenting path $P$ from a free point $b\in B_\cell$ to a point $b'\in B_\cell$, the net-cost of $P$ is 
\[\phi(P)=y(b) + s(b') +\sum_{(a'',b'')\in P}s(a'',b'') \ge \kappa_{b'} + s(b'),\]
and we get an equality if $P$ is a shortest path from $s$ to $b'$.
Therefore, the updated key of $\cell$, i.e., $\kappa_\cell=\min\{\min_{a\in A_\cell^F}\kappa_a, \min_{b\in B_\cell}\kappa_b+s(b)\}$ correctly computes the net-cost of the minimum net-cost augmenting path inside $\cell$. Finally, note that for any augmenting path $P$ from a free point $b\in B_\cell$, by Lemma~\ref{lem:slackcost}, $\phi(P)\ge y(b)=\ymax$ and the key of $\cell$ would be at least $\ymax$.









    

% Property (LP3) shows that the merge step on $\cell$ terminates. Suppose $M_\cell, y(\cdot)$ denotes the matching initially constructed by the merge step by combining the matchings maintained inside the children of $\cell$. By Lemma~\ref{lemma:combination}, the matching $M_\cell, y(\cdot)$ is $\cell$-feasible. Furthermore, assuming invariant (I2) before the execution of the merge step on $\cell$, all points of $B_\cell$ have dual weights at most $\ymax$. Finally, by invariant (I1) for the children of $\cell$, $y(a)=0$ for all free points $a\in \freeofcell{A}{\cell}$. Hence, the matching $M_\cell, y(\cdot)$ satisfies the conditions of Lemma~\ref{lemma:local}. 

% By applying Lemma~\ref{lemma:local} to each iteration of the merge step on $\cell$, when the merge step terminates, the matching $M_\cell, y(\cdot)$ is a $\cell$-feasible matching (from (LP1)) where $y(b)\le\ymax$ for all points $b\in B_\cell$ and $y(a)=0$ for all free points $a\in \freeofcell{A}{\cell}$ (from (LP2)). Furthermore, by the termination condition, any active free point $b\in B_\cell$ has a dual weight $y(b)=\ymax$. Therefore, the matching $M_\cell, y(\cdot)$ is a $\cell$-optimal matching and invariants (I1) and (I2) continue to be true.  

% For any cell $\cell$, let $M_{\cell'}, y(\cdot)$ denote the $\cell'$-optimal matching computed inside each child $\cell'$ of $\cell$ before executing the merge step on $\cell$. In the following, we show that
% \begin{itemize}[itemsep=2pt,parsep=0pt,topsep=0pt,leftmargin=5em]
%     \item[(LP1)] the matching $M_\cell, y(\cdot)$ is a $\cell$-optimal matching, and
%     \item[(LP2)] $y(b)=\ymax$ for all active free points $b\in B_\cell$. 
% \end{itemize}

% Next, we show that invariants (I1) and (I2) continue to be true after each execution of the extended Hungarian search procedure.






% % Therefore, invariants (I1) and (I2) remain true after each execution of the extended Hungarian search procedure.

% \paragraph{Proof of Invariant (I3) and Lemma~\ref{lemma:phi-value}.} 
% We use an inductive argument to prove both invariants (I3) and also show that the value $\ymax$ maintained by our algorithm as the minimum key value of the cells in $\pq$ is non-decreasing.

% For any cell $\cell$ of $\mcH$, the termination condition of the merge step on $\cell$ guarantees that all active free points $b\in B_\cell$ have $y(b)=\ymax$. By Lemma~\ref{lem:slackcost}, for the $\cell$-optimal matching $M_\cell, y(\cdot)$ computed inside $\cell$ and any $\cell$-augmenting path $P$ from an active free point $b\in B_\cell$, $\phi(P)\ge y(b)=\ymax$; therefore, the key of $\cell$ in $\pq$ after locally processing $\cell$ will be at least $\ymax$, and the merge step does not result in a modification to $\ymax$. Additionally, for all other current cells $\cell'\neq\cell$, assuming (I3) prior to locally processing $\cell$ and since $\ymax$ has not changed, it will hold after the merge step as well. 

% Next, we show the same properties for the extended Hungarian search procedure. For this step, our algorithm extracts a cell $\cell$ with the minimum key value, which is our $\ymax$ value, and globally processes $\cell$. Upon executing the extended Hungarian search procedure on $\cell$, by Lemma~\ref{lemma:global}, all active free points $b\in B_\cell$ will have $y(b)=\ymax$. The algorithm then augments $M_\cell$ along an admissible $\cell$-augmenting path and updates the key of $\cell$ in $\pq$. By Lemma~\ref{lem:slackcost}, the net-cost of the next minimum net-cost $\cell$-augmenting path inside $\cell$ is lower-bounded by the dual weight of the active free points, which is $\ymax$; hence, the new key value of $\cell$ in $\pq$ is at least $\ymax$ and the extended Hungarian search procedure does not decrease $\ymax$. 

% To complete our proof for invariant (I3), it remains to show that for any current cell $\cell'\neq \cell$, the invariant (I3) for $\cell'$ remains true after the extended Hungarian search procedure on $\cell$.
% % On the other hand, the extended Hungarian search procedure might change the value of $\ymax$, since it requires extracting the minimum value of $\pq$. 
% % Suppose $\cell'\neq\cell$ is the extracted cell. Assuming that invariant (I3) holds prior to globally processing $\cell'$, there exists a set of dual weights $y'(\cdot)$ where $M_{\cell'}, y'(\cdot)$ is $\cell'$-optimal and $y'(b)=\ymax$ for each active free point $b\in B$, where $\ymax$ denotes the value maintained before globally processing $\cell'$. Since the net-cost of any $\cell'$-augmenting path is at least as high as the dual weight of its active free endpoint, for the minimum net-cost path $P_{\cell'}$, $\phi(P_{\cell'})\ge \ymax$, and therefore, by globally processing $\cell'$, the value of $\ymax$ does not decrease.
% Let $\ymax_{\cell'}$ denote the net-cost of the minimum net-cost $\cell'$-augmenting path inside $\cell'$. Since $\cell$ is extracted from $\pq$ with a key $\ymax$, $\ymax_{\cell'}\ge \ymax$. In the following, we construct a set of the dual weights $y'(\cdot)$ so that $M_{\cell'}, y'(\cdot)$ is $\cell'$-optimal and the dual weight of the active free points of $B_{\cell'}$ become $\ymax$. 

% Let $\kappa_u$ denote the distance of each point $u\in A_{\cell'}\cup B_{\cell'}$ from the source vertex $s$ in the residual graph. Define $y'(u)\leftarrow y(u) + \ymax - \kappa_u$. This completes the construction of the dual weights. Since $\ymax\le \ymax_{\cell'}=\min\{\min_{a\in\freeofcell{A}{\cell'}}\kappa_a, \min_{b\in B_{\cell'}}\kappa_b+s_{\cell'}(b)\}$, the dual updates do not change the dual weight of any free point $a\in \freeofcell{A}{\cell'}$ and also no point $b\in B_{\cell'}$ can get a dual weight more than $\distance{b}{\cell'}$; hence, using a proof similar to the proof of Lemma~\ref{lemma:global}, one can show that the matching $M_{\cell'}, y'(\cdot)$ is $\cell'$-optimal. Furthermore, since $\kappa_b=y(b)$ for all active free points $b\in B_{\cell'}$, $y'(b)=y(b) + \ymax - \kappa_b=\ymax$, as claimed.

% \begin{restatable}{lemma}{phiValue}\label{lemma:phi-value}
% The value $\ymax$ maintained by our algorithm is non-decreasing.
% \end{restatable}
% \begin{proof}
%     From Lemma~\ref{lemma:hung_properties}, for the cell $\cell_{\min}$ with the smallest key value in $\pq$, after executing the extended Hungarian search procedure on $\cell_{\min}$, its updated key value is at least $\varphi$. Furthermore, from Lemma~\ref{lemma:merge_properties}, after the merge step on any cell $\cell'$, the key value of $\cell'$ in $\pq$ (if $\cell'$ is inserted to $\pq$) is at least $\varphi$. Hence, the value of $\varphi$ is non-decreasing during the execution of our algorithm.
% \end{proof}




% % \begin{restatable}{lemma}{localProp}\label{lemma:local}
% %     For any cell $\cell$, let $M_\cell, y(\cdot)$ denote the matching computed inside $\cell$ after locally processing $\cell$. Then,
% % \begin{itemize}[itemsep=2pt,parsep=0pt,topsep=0pt,leftmargin=5em]
% %     \item[(LP1)] the matching $M_\cell, y(\cdot)$ is a $\cell$-optimal matching, and
% %     \item[(LP2)] $y(b)=\ymax$ for all active free points $b\in B_\cell$. 
% % \end{itemize}
% % \end{restatable}
% % \begin{proof}
    
% \end{proof}

\subsection{Runtime Analysis of the Merge Step} 
For any non-leaf cell $\cell$ with $\cell'$ and $\cell''$ as children, the merge step at $\cell$ first increases the dual weights of the free points inside $\cell'$ and $\cell''$ in $\tilde{O}(n_\cell\Phi(n_\cell))$ time (Lemma~\ref{lemma:fresh_duals}). For the feasible extended matching $M^\mcC=(M, B^\mcC), y(\cdot)$ after this initial step, let $B^\mcC_\cell$ denote the subset of the boundary-matched points in $B^\mcC$ that are matched to $\Gamma_\cell$. As discussed in Section~\ref{subsec:merge_correctness}, each iteration of the while-loop in the merge step reduces the number of free points with a dual weight less than $\ymax$ by one; therefore, the total number of executions of the while-loop in the merge step is at most $|B^\mcC_\cell|$.

\begin{restatable}{lemma}{localProcessInactive}\label{lemma:local-inactive}
    For any cell $\cell$, the number of iterations of the merge step on $\cell$ is $O(|B^\mcC_\cell|)$.
\end{restatable}

Each iteration requires the computation of the distance $\kappa_v$ for each point $v\in A_\cell\cup B_\cell$, which takes $\tilde{O}(n_\cell^2)$ time. As shown in Section~\ref{sec:hungarian_nk}, the efficiency of this computation can be improved to $\tilde{O}(n_\cell\Phi(n_\cell))$ using a dynamic weighted nearest neighbor data structure with a query/update time of $\Phi(n_\cell)$. Furthermore, by Lemma~\ref{lemma:combination}, $|B_{\cell}^\mcC|=O(n^{4/5})$. Computing the updated key of $\cell$ also requires the computation of the distance of each point from the source in the residual graph, which also takes $\tilde{O}(n_\cell\Phi(n_\cell))$ time. Therefore, the total execution time of the merge step would be $\tilde{O}(n^{4/5}n_\cell\Phi(n_\cell))$.



% For any cell $\cell$ of $\mcH$, let $M_{\cell'}, y(\cdot)$ denote the $\cell'$-optimal matching computed inside each child $\cell'$ of $\cell$. Let $B^\mcC_\cell$ denote the set of active free points with respect to $M_{\cell'}, y(\cdot)$ for each child $\cell'$ of $\cell$. Let $M_\cell=\bigcup_{\cell'\in\children\cell}M_{\cell'}$ denote the matching computed at the beginning of the merge step on $\cell$. Recall that $\divider\cell$ denotes the divider of $\cell$. For any child $\cell'\in\children\cell$, let $\mcB^{\divider\cell}_{\cell'}$ denote the set of inactive free points of $B_{\cell'}$ that are mapped to $\divider\cell$ with respect to $M_{\cell'}, y(\cdot)$. Define $m_\cell=\sum_{\cell'\in\children\cell}|\mcB^{\divider\cell}_{\cell'}|$ as the total number of inactive free points inside the children of $\cell$ that are mapped to the divider $\divider\cell$ of $\cell$. 

% The matching $M_\cell, y(\cdot)$ is a $\cell$-feasible matching, and $B^\mcC_\cell\cup\bigcup_{\cell'\in\children\cell}\mcB_{\cell'}^{\divider\cell}$ is the set of all active free points with respect to $M_\cell, y(\cdot)$. Note that for any child $\cell'$ of $\cell$, the net-cost of the minimum net-cost $\cell'$-augmenting path inside $\cell$ is at least $\ymax$; therefore, the active free points with respect to $M_{\cell'}, y(\cdot)$ need to increase at least as high as $\ymax$ to find an admissible $\cell'$-augmenting path. In this case, no point in the set $B^\mcC_\cell$ will form an admissible $\cell$-augmenting path during the merge step and the only points that might form admissible $\cell$-augmenting paths are the ones in the set $\mbB_\cell:=\bigcup_{\cell'\in\children\cell}\mcB_{\cell'}^{\divider\cell}$, leading to the following lemma.



% To bound the running time of locally processing a cell $\cell$, we show in the following that $m_\cell=O(n^{4/5})$. To do so, we first show in the following lemma that there exists a matching $M$ on $\mcG_\cell$ that matches the majority of points and has a low cost. \todo{fix}





% To bound the number of points in $B_{\mathrm{aug}}$, we show in Equation~\eqref{eq:net-cost-sum-2} below that the total dual weights of the points in $B_{\mathrm{aug}}$ is at most $w(M')$. 
% Using this observation, we show in Equations~\eqref{eq:close-points} and~\eqref{eq:far-points} that the number of points in $F_{\mathrm{aug}}(M_\cell)$ is $\tilde{O}(n^{4/5})$. Combining the two bounds, $|F(M_\cell)| = \tilde{O}(n^{4/5})$, as claimed. 


% {\it (i) Matching Construction.} We begin by constructing a partial matching $M$ that has a high cardinality and a low cost.


% To prove this lemma, we construct a matching $M$ as follows. Consider a sequence of grids $\langle G_1, \ldots, G_t\rangle$ ordered in increasing order of their cell side length. Starting from the smallest grid $G_1$, we arbitrarily match as many points as possible inside each square of $G_1$. We then move to the second grid $G_2$ and do the same, i.e., for each square of $G_2$, we arbitrarily match as many unmatched points as possible inside the square. We continue this process until we process grid $G_t$. We show that by carefully picking the cell side length of the grids, the matching $M$ will match all except $O(n^{3/4})$ points with a cost $O(\ell_\cell^2 n^{1/4})$. We provide full proof in Appendix~\ref{sec:appendix-GRS}.

% {\it (ii) Bounding the Total Dual Weights.}
% We next show that the total dual weights of the free points in $B_{\mathrm{aug}}$ is at most $w(M')$. To do so, we begin by showcasing an important property of $\cell$-optimal matchings, whose proof is included in the appendix.
% \begin{itemize}
%     \item[(P1)] For any augmenting path $P$ from a free point $b\in B_\cell$ to a point $a\in A_\cell$, $\phi(P)\ge y(b)$.
% \end{itemize}
% Let $\mcP$ denote the set of all augmenting paths in the symmetric difference $M\oplus M_\cell$, i.e., $\mcP$ consists of all paths $P=\langle b_1, a_1, b_2, \ldots, b_t, a_t\rangle$ such that for all $i\in[1,t]$, $(b_i, a_i)\in M$ and for all $i\in[1, t-1]$, $(a_i, b_{i+1})\in M_\cell$. 
% % Note that $|F(M_\cell)\setminus F(M)|=|\mcP|$. 
% From the definition of net-costs,
% \begin{align}
%     \sum_{p\in\mcP} \phi(P) &= \sum_{p\in\mcP} \left(\sum_{(a,b)\in P\cap M}\distance{a}{b} - \sum_{(a,b)\in P\cap M_\cell}\distance{a}{b} \right)\le \sum_{p\in\mcP} \left(\sum_{(a,b)\in P\cap M}\distance{a}{b}\right)\le w(M).\label{eq:net-cost-sum-1}
% \end{align}
% For each $P\in\mcP$, let $b_P\in B_\cell$ and $a_P\in A_\cell$ denote the two end-points of $P$. Using property (P1) and Equation~\eqref{eq:net-cost-sum-1}, 
% % Recall that $F(M)$ (resp $F(M_\cell)$) denotes the set of all free points of $B_\cell$ with respect to $M$ (resp. $M_\cell$). 
% \begin{align}
%     w(M)\ge \sum_{P\in\mcP} \phi(P) &\ge \sum_{P\in\mcP} y(b_P)= \sum_{b\in F_{\mathrm{aug}}(M_\cell)} y(b).\label{eq:net-cost-sum-2}
% \end{align}

% {\it (iii) Bounding the Number of Iterations.}
% Define $\theta:=\ell_\cell n^{-1/5}$. Each free point $b\in F(M_\cell)$ is called a \emph{close} (resp. \emph{far}) point if the Euclidean distance of $b$ to the divider of $\cell$ is at most (resp. more than) $\theta$. Let $F_{\mathrm{aug}}^{\mathrm{close}}(M_\cell)$ (resp. $F_{\mathrm{aug}}^{\mathrm{far}}(M_\cell)$) denote the set of all close (resp. far) points of $F_{\mathrm{aug}}(M_\cell)$. By the construction of the hierarchical partitioning,
% \begin{equation}\label{eq:close-points}
%     |F_{\mathrm{aug}}^{\mathrm{close}}(M_\cell)| = O(n_\cell n^{-1/5}) = O(n^{4/5}).
% \end{equation}
% For each far point $b\in F_{\mathrm{aug}}(M_\cell)$, since each active free point in $B_\cell$ is mapped to the divider $\divider\cell$ in one child of $\cell$, $y(b)=\distance{b}{\divider\cell}\ge \theta$; therefore, 
% \begin{equation}\label{eq:net-cost-3}
%     \sum_{b\in F_{\mathrm{aug}}(M_\cell)} y(b) \ge \sum_{b\in Y} y(b) \ge \theta\times |F_{\mathrm{aug}}^{\mathrm{far}}(M_\cell)|.
% \end{equation}
% Combining Equations~\eqref{eq:net-cost-sum-2} and~\eqref{eq:net-cost-3},
% \begin{equation}\label{eq:far-points}
%     |F_{\mathrm{aug}}^{\mathrm{far}}(M_\cell)| \le \frac{\sum_{b\in F_{\mathrm{aug}}(M_\cell)} y(b)}{\theta} \le \frac{w(M)}{\theta} = O(n^{4/5}).
% \end{equation}
% Combining with Equation~\eqref{eq:close-points},
% \begin{align*}
%     |F(M_\cell)| &\le |F_{\mathrm{aug}}(M_\cell)| + |F_{\mathrm{alt}}(M_\cell)| \le |F_{\mathrm{aug}}^{\mathrm{close}}(M_\cell)| + |F_{\mathrm{aug}}^{\mathrm{far}}(M_\cell)| + |F(M)| = \tilde{O}(n^{4/5}).
% \end{align*}


\subsection{Runtime Analysis of the Extended Hungarian Search Procedure}\label{appendix:hung_time}
Recall that the extended Hungarian search procedure iteratively picks the cell with the minimum key from $\pq$ to be processed.
For each leaf cell $\cell$ in $\mcH$, since $\cell$ contains the points corresponding to a single request and contains only one point of $B$, the procedure picks $\cell$ at most once, at which it matches the point $b\in B_\cell$ to the boundaries of $\cell$. Therefore, the total time of the extended Hungarian search procedure on all leaf cells of $\mcH$ is $O(n)$. 

Next, we show that for any non-leaf cell $\cell$ of $\mcH$, our algorithm executes the search procedure on $\cell$ at most $O(n^{4/5})$ times. Since each execution of the procedure takes $\tilde{O}(n_{\cell}\Phi(n))$ time and each point participates in $O(\log (n\Delta))$ cells, the total running time of the extended Hungarian search procedure on all cells of $\mcH$ would be $O(n^{4/5}\sum_{\cell\in\mcH}n_\cell\Phi(n_\cell)) = O(n^{9/5}\Phi(n_\cell)\log(n\Delta))$, as claimed.

For any execution of the search procedure on $\cell$, since $\cell$ has the minimum key in $\pq$, the value $\ymax$ represents the net-cost of the minimum net-cost augmenting path inside $\cell$.
Recall that for any non-leaf cell $\cell$, we categorized the selections of $\cell$ by the search procedure as low-net-cost if the value of $\ymax$ in that iteration is at most $\ell_\cell n^{-1/5}$ and high-net-cost otherwise. 
To bound the high-net-cost selections, we show that as soon as the value $\ymax$ exceeds $\ell_\cell n^{-1/5}$, the number of free points inside $\cell$ cannot be more than $O(n^{4/5})$, and therefore, the number of high-net-cost selections of $\cell$ is $O(n^{4/5})$.

\highMax*
\begin{proof}
    Let $M'$ be the matching inside $\cell$ as constructed in Lemma~\ref{lemma:geometric-matching}, and let $M$ denote the matching of the extended matching $M^\mcC=(M, B^\mcC)$.
    Define $B^\mcC_\cell$ as the set of free points of $B_\cell$ with respect to $M^\mcC,y(\cdot)$. Let $M_\cell$ denote the subset of matching edges of $M$ that lie inside $\cell$, and let $\mcP_{\mathrm{aug}}$ (resp. $\mcP_{\mathrm{alt}}$) denote the set of (standard) augmenting (resp. alternating) paths in the symmetric difference $M_\cell\oplus M'$ with one endpoint in the set $B^\mcC_\cell$. Note that each path in $\mcP_{\mathrm{alt}}$ has one endpoint that is free in $M'$ and therefore, $|\mcP_{\mathrm{alt}}|\le |F(M')|=O(n^{4/5})$; here, $F(M')$ denotes the set of free points of $M'$. Next, we show that $|\mcP_{\mathrm{aug}}|=O(n^{4/5})$.

    From the definition of the net-cost of an augmenting path,
    \begin{align}
        \sum_{P\in \mcP_{\mathrm{aug}}} \phi(P) &= \sum_{P\in \mcP_{\mathrm{aug}}} \left(\sum_{(a,b)\in P\cap M'}\distance{a}{b} - \sum_{(a,b)\in P\cap M_\cell}\distance{a}{b} \right)\nonumber\\ &\le \sum_{P\in \mcP_{\mathrm{aug}}} \left(\sum_{(a,b)\in P\cap M'}\distance{a}{b}\right)\le w(M').\label{eq:net-cost-sum-1-1}
    \end{align}
    For each path $P\in \mcP_{\mathrm{aug}}$, let $b_P\in B_\cell$ and $a_P\in A_\cell$ denote the two end-points of $P$. Define $B_{\mathrm{aug}}:=\{b_P:P\in \mcP_{\mathrm{aug}}\}$ as the set of free endpoints of the paths in $\mcP_{\mathrm{aug}}$. Using Lemma~\ref{lem:slackcost} and Equation~\eqref{eq:net-cost-sum-1-1}, 
    \begin{align}
        \sum_{b\in B_{\mathrm{aug}}} y(b) = \sum_{P\in \mcP_{\mathrm{aug}}} y(b_P)\le \sum_{P\in \mcP_{\mathrm{aug}}} \phi(P) \le w(M') .\label{eq:net-cost-sum-2-1}
    \end{align}
    From invariant (I2), the dual weight of all free points of $B_\cell$ equals $\ymax>\ell_\cell n^{-1/5}$. Therefore,
    \begin{equation*}\label{eq:net-cost-sum-2-2}
        |B_{\mathrm{aug}}| = \frac{\sum_{b\in B_{\mathrm{aug}}} y(b)}{\ymax} \le \frac{w(M')}{\ell_\cell n^{-1/5}}=O(n_\cell^{3/5}n^{1/5}) = O(n^{4/5}).
    \end{equation*}
    Combining the two bounds, the total number of free points with respect to $M^\mcC, y(\cdot)$ is $|\mcP_{\mathrm{alt}}|+|\mcP_{\mathrm{aug}}|=O(n^{4/5})$.
\end{proof}

Define $\mcC_\cell$ as the set of all cells of $\mcH$, including $\cell$ itself, that are processed by the merge step of our algorithm while $\cell\in\mcC$ and $\ymax\le \ell_\cell n^{-1/5}$. For any cell $\cell'\in\mcC_\cell$, let $\mcB_{\cell'}$ denote the set of points of $B$ that are matched to the divider $\Gamma_{\cell'}$ at the beginning of the merge step at $\cell'$. During the execution of our algorithm, before processing $\cell'$ by the merge step, there are $k$ free points across all cells is $\mcC$. After combining the two children of $\cell'$ (and removing the divider of $\cell'$), the number of free points across all cells in the partitioning is now increased to $k + |\mcB_{\cell'}|$. Each iteration of the merge step either (i) finds an admissible augmenting path, which reduces the number of free points by one, or (ii) finds an admissible alternating path to a point $b\in B_{\cell'}$ with $y(b)=\ymax$, which does not change the number of free points. Therefore, after the merge step at $\cell'$, the number of free points across all cells in the partitioning is at most $k + |\mcB_{\cell'}|$. Our algorithm then iteratively executes the search procedure to reduce the number of free points across all cells to $k$. Hence, for each cell $\cell'\in\mcC_\cell$, the merge step at $\cell'$ might lead to the execution of a low-net-cost extended Hungarian search procedure on $\cell$ at most $|\mcB_{\cell'}|$ times.

\begin{restatable}{lemma}{iterationsOfGlobal}\label{lemma:global-iterations}
    For any cell $\cell$, the number of low-net-cost executions of the extended Hungarian search procedure on $\cell$ is at most $\sum_{\cell'\in\mcC_\cell} |\mcB_{\cell'}|$.
\end{restatable}

Recall that in a low-net-cost selection of $\cell$, the value $\ymax\le \ell_\cell n^{-1/5}$. Using invariant (I2), the dual weight of all points of $B$ is at most $\ymax$, and therefore, for any cell $\cell'\in\mcC_{\cell}$ and each point $b\in \mcB_{\cell'}$, $y(b)\le \ell_\cell n^{-1/5}$, i.e., $b$ is a boundary-matched point that is matched to the divider $\divider{\cell'}$ and has a dual weight at most $\ell_\cell n^{-1/5}$. Therefore, $\distance{b}{\divider{\cell'}} = y(b) \le \ell_\cell n^{-1/5}$ and all points in $\mcB_{\cell'}$ are at a distance at most $\ell_\cell n^{-1/5}$ from the divider $\divider{\cell'}$. From Lemma~\ref{lemma:ratio}, $\ell_{\cell}\le 3\ell_{\cell'}$. Therefore, using Lemma~\ref{lemma:margin}, the number of points of $B_{\cell'}$ at a distance at most $\lambda\ell_{\cell'}\ge \ell_\cell n^{-1/5}$ to the divider $\divider{\cell'}$ is $O(n_{\cell'}n^{-1/5})$; hence, $|\mcB_{\cell'}| = O(n_{\cell'}n^{-1/5})$.

 
\ratio*
% \begin{proof}
%     For the cell $\cell$ (resp. $\cell'$), suppose $\cell_{\min}$ (resp. $\cell'_{\min}$) denote the smaller child of $\cell$ (resp. $\cell'$). Since our algorithm picked $\cell_{\min}$ for being merged as the smallest cell before $\cell'_{\min}$, 
%     \begin{equation*}\label{eq:ratio-1}
%         \ell_\cell\le 3\ell_{\cell_{\min}}\le 3\ell_{\cell'_{\min}}\le 3\ell_{\cell'},
%     \end{equation*}
%     where the first and third inequalities follow by the construction of our hierarchical partitioning.
%     Similarly, 
%     \begin{equation*}\label{eq:ratio-2}
%         \ell_{\cell'}\le 3\ell_{\cell'_{\min}}\le 3\ell_{\cell},
%     \end{equation*}
%     where the first inequality holds by the construction of $\mcH$ and the second inequality holds since our algorithm picked $\cell'_{\min}$ as the smallest cell to be processed rather than $\cell$.
% \end{proof}
\begin{proof}
    To prove this lemma, we first provide two useful relations between the sum of side-lengths $p_{\hat{\cell}}$ and the largest side-length $\ell_{\hat{\cell}}$ of any cell $\hat{\cell}$ of $\mcH$. Let $\ell_x$ (resp. $\ell_y$) denote the width (resp. height) of $\hat\cell$. Since the aspect ratio of $\hat{\cell}$ is at most $3$, i.e., $\min\{\ell_x, \ell_y\} \ge \frac{1}{3}\ell_{\hat{\cell}} = \frac{1}{3}\max\{\ell_x, \ell_y\}$, 
    \begin{equation*}\label{eq:ratio-edge}
        \frac{4}{3}\ell_{\hat{\cell}} \le \ell_x + \ell_y = p_{\hat{\cell}} =\ell_x + \ell_y \le 2\ell_{\hat{\cell}}.
    \end{equation*}
    Furthermore, for the smaller child $\hat{\cell}_1$ of $\hat{\cell}$, if the cell is divided on the $x$ axis, then the width of $\hat{\cell}_1$ is within  $\frac{1}{3}\ell_{\hat\cell}$ and $\frac{1}{2}\ell_{\hat\cell}$; therefore,
    \begin{equation*}
        \frac{4}{3}p_{\hat{\cell}_1}\le p_{\hat\cell} \le 2 p_{\hat{\cell}_1}.
    \end{equation*}
    For the cell $\cell$ (resp. $\cell'$), suppose $\cell_{\min}$ (resp. $\cell'_{\min}$) denote the smaller child of $\cell$ (resp. $\cell'$). Since our algorithm picked $\cell_{\min}$ for being merged as the smallest cell before $\cell'_{\min}$, 
    \begin{equation}\label{eq:ratio-1}
        \frac{4}{3}\ell_\cell\le p_\cell\le 2p_{\cell_{\min}}\le 2p_{\cell'_{\min}}\le \frac{3}{2}p_{\cell'}\le 3\ell_{\cell'}.
    \end{equation}
    Similarly, since our algorithm picked $\cell'_{\min}$ as the smallest cell to be processed rather than $\cell$,
    \begin{equation}\label{eq:ratio-2}
        \frac{4}{3}\ell_{\cell'}\le p_{\cell'}\le 2p_{\cell'_{\min}}\le 2p_{\cell}\le 4\ell_\cell.
    \end{equation}
    Combining Equations~\eqref{eq:ratio-1} and~\eqref{eq:ratio-2},
    \begin{equation*}
        \frac{1}{3} \ell_{\cell'} \le \ell_\cell \le \frac{9}{4} \ell_{\cell'}.
    \end{equation*}
\end{proof}

% Since each cell of $\mcH$ has a constant aspect ratio, we get the following lemma.



We next show that $\sum_{\cell'\in\mcC_{\cell}}n_{\cell'}=O(n)$ and conclude that $\sum_{\cell'\in\mcC_\cell}|\mcB_{\cell'}| = O(n^{4/5})$, which bounds the number of low-net-cost executions of the search procedure on $\cell$ by $O(n^{4/5})$, as desired.
By Lemma~\ref{lemma:ratio}, for any cell $\cell'\in\mcC_{\cell}$, $\ell_{\cell'}\in[\frac{1}{3}\ell_\cell, \frac{9}{4}\ell_\cell]$. By the construction of $\mcH$, for any cell $\cell'$ and its grandparent $\cell'_g$, $\ell_{\cell'}\le \frac{2}{3}\ell_{\cell'_g}$. Therefore, for $\cell'\in \mcC_\cell$, only the ancestor of $\cell'$ up to $2\log_{3/2} \frac{27}{4}$ levels higher can also be in $\mcC_\cell$; therefore, for any point $u\in A\cup B$, the point $u$ lies inside at most $O(1)$ cells of $\mcC_\cell$, and $\sum_{\cell'\in\mcC_{\cell}}n_{\cell'}=O(n)$, as claimed.

We conclude that the total number of executions of the extended Hungarian search for each cell $\cell$ of $\mcH$ is $O(n^{4/5})$. The next lemma follows since each execution takes $\tilde{O}(n_\cell\Phi(n))$ time.
\begin{restatable}{lemma}{GlobalTime}\label{lemma:global-time}
    For any cell $\cell$ of $\mcH$, the total execution time of the extended Hungarian search procedure on $\cell$ takes $\tilde{O}(n^{4/5}n_\cell\Phi(n))$ time.
\end{restatable}











\subsection{Extension to higher dimensions}\label{sec:appendix-high-d-ksp}
Given a set of requests $\requests$ in the $d$ dimensions, for any $d\ge 2$, let $\mcG_\requests$ on point sets $A$ and $B$ denote the graph constructed for the $k$-SP problem on $\requests$ under any $\ell_p$ norm for some $p\ge 1$. Let $\mcH_d$ denote the hierarchical partitioning constructed for the point set $A\cup B$ with a parameter $\lambda = 9n^{-\frac{1}{2d+1}}$. 
The hierarchical partitioning has a height $O(d\log(n\Delta))$. Using $\mcH_d$, we run our algorithm as described in Section~\ref{sec:k-seq}. 

We summarize the efficiency analysis of our algorithm for $d$-dimensional point sets next and show that our algorithm runs in $\tilde{O}(n^{2-\frac{1}{2d+1}}\Phi(n)\log \Delta)$ time. In particular, we show that for any cell $\cell$, the number of iterations of the merge step on $\cell$ is $O(n^{1-\frac{1}{2d+1}})$, where each iteration takes $\tilde{O}(n_\cell\Phi(n))$ time. We also show that our algorithm runs the extended Hungarian search procedure in $O(n^{1-\frac{1}{2d+1}})$ time on $\cell$, each in $\tilde{O}(n_\cell\Phi(n))$ time. Adding these bounds for all cells of $\mcH_d$, the total running time of our algorithm is $\tilde{O}(n^{2-\frac{1}{2d+1}}\Phi(n)\log\Delta)$. We summarize the details below.

Given a feasible extended matching $M^\mcC, y(\cdot)$, for any cell $\cell\in \mcC$, let $B^\mcC_\cell$ denote the set of all points of $B_\cell$ that are matched to the divider $\divider\cell$ in the matching before the execution of the merge step at $\cell$. By Lemma~\ref{lemma:local-inactive}, the number of iterations of the merge step on $\cell$ is $O(|B^\mcC_\cell|)$.
To bound the number of points in $B^\mcC_\cell$, we first show that there exists a partial matching $M'$ of high cardinality and low cost. 
\begin{lemma}\label{lemma:geometric-matching-d}
    For any cell $\cell$ of $\mcH_d$, there exists a matching $M'$ over $\mcG_\sigma$ inside $\cell$ that matches all except $O(n_\cell^{1-\frac{1}{2d+1}})$ points of $B_\cell$ and has a cost $O(\ell_\cell n_\cell^{1-\frac{2}{2d+1}})$.
\end{lemma}
\begin{proof}
    Similar to our construction for Lemma~\ref{lemma:geometric-matching}, to construct the matching $M'$ for Lemma~\ref{lemma:geometric-matching-d}, we place a grid of cell-side-length $\ell_\cell n_\cell^{-\frac{2}{2d+1}}$ and compute the matching corresponding to the $1$-partitioning of the requests inside each cell of this grid. It is easy to confirm that the matching $M'$ achieves the bounds claimed in Lemma~\ref{lemma:geometric-matching-d}.
\end{proof}


For the feasible extended matching $M^\mcC=(M, B^\mcC), y(\cdot)$, let $M_\cell$ denote the subset of matching edges of $M$ that lie inside $\cell$. 
Let $\mcP_{\mathrm{aug}}$ (resp. $\mcP_{\mathrm{alt}}$) denote the set of augmenting paths (resp. alternating paths) with an endpoint in $|B^\mcC_\cell|$ in the symmetric difference $M_\cell\oplus M'$.
As discussed above, $|\mcP_{\mathrm{alt}}|\le |F(M')|=O(n^{1-\frac{1}{2d+1}})$, where $F(M')$ denotes the set of free points of $B_\cell$ with respect to $M'$. 
We partition the free endpoints of $\mcP_{\mathrm{aug}}$ into the set $B^{\mathrm{close}}_{\mathrm{aug}}$ (resp. $B^{\mathrm{far}}_{\mathrm{aug}}$) that are at a distance closer than (resp. further than) $\lambda'_\cell=\ell_\cell n^{-\frac{1}{2d+1}}$ to the divider $\divider\cell$ of $\cell$. From Lemma~\ref{lemma:margin},  $|B^{\mathrm{close}}_{\mathrm{aug}}|=O(n^{1-\frac{1}{2d+1}})$. Finally, by Equation~\eqref{eq:net-cost-sum-2}, $\sum_{b\in B^{\mathrm{far}}_{\mathrm{aug}}}y(b) \le w(M')$. Since the dual weight of each free point in $B^{\mathrm{far}}_{\mathrm{aug}}$ is at least $\lambda'_\cell$, we get a bound of $O(n^{1-\frac{1}{2d+1}})$ on the number of such points, and a total execution time of $\tilde{O}(n^{1-\frac{1}{2d+1}}n_\cell\Phi(n))$ for the merge step on a cell $\cell$ of $\mcH$.




Next, we bound the number of executions of the extended Hungarian search procedure for each cell $\cell$.
A selection of $\cell$ by the extended Hungarian search procedure is low-net-cost if $\ymax\le \ell_\cell n^{-\frac{1}{2d+1}}$ and high-net-cost otherwise. For the high-net-cost selections, one can use the matching $M'$ from Lemma~\ref{lemma:geometric-matching-d} and a similar argument as in Lemma~\ref{lemma:high-max} to show that when $\ymax> \ell_\cell n^{-\frac{1}{2d+1}}$, the total number of free points inside $\cell$ cannot exceed $O(n^{1-\frac{1}{2d+1}})$ and conclude an upper bound of $O(n^{1-\frac{1}{2d+1}})$ on the number of high-net-cost selections of $\cell$. 
% \begin{restatable}{lemma}{highMax}\label{lemma:high-max}
%     For any cell $\cell\in\mcH$, let $M_\cell, y(\cdot)$ denote a $\cell$-optimal matching on $\mcG_\cell$. If $\max_{b\in B_\cell}y(b)>\ell_\cell n^{-1/5}$, then the number of active free points of $M_\cell,y(\cdot)$ is $O(n^{4/5})$.
% \end{restatable}
To bound the number of low-net-cost selections of $\cell$, define $\mcC_\cell$ as the set of all cells $\cell'\in\mcH$ that are processed by the merge step while $\cell$ is in $\mcC$ and $\ymax\le \ell_\cell n^{-\frac{1}{2d+1}}$. For any cell $\cell'\in\mcC_\cell$, let $\mcB_{\cell'}$ denote the set of boundary-matched points of $B_{\cell'}$ matched to the divider of $\cell'$ before the execution of the merge step on $\cell'$. By Lemma~\ref{lemma:global-iterations}, the number of low-net-cost selections of $\cell$ is at most $\sum_{\cell'\in\mcC_\cell}|\mcB_{\cell'}|$. Since $\ymax\le \ell_\cell n^{-\frac{1}{2d+1}}$, for each cell $\cell'\in\mcC_\cell$, all points in $\mcB_{\cell'}$ are at a distance at most $\ell_\cell n^{-\frac{1}{2d+1}}$ to the divider of $\cell'$. Therefore, from Lemma~\ref{lemma:margin} and~\ref{lemma:ratio} and using an identical discussion as above, $\sum_{\cell'\in\mcC_\cell}\mcB_{\cell'} = O(dn^{1-\frac{1}{2d+1}})$. Therefore, the total execution time of the extended Hungarian search step on any cell $\cell$ of $\mcH$ is $\tilde{O}(dn^{1-\frac{1}{2d+1}}n_\cell\Phi(n))$. 

Combining the total execution times of the merge step and the extended Hungarian search step, the running time of our algorithm would be $\tilde{O}(\sum_{\cell\in\mcH}dn^{1-\frac{1}{2d+1}}n_\cell\Phi(n)) = \tilde{O}(dn^{2-\frac{1}{2d+1}}\Phi(n)\log\Delta)$, leading to the following theorem.

\begin{theorem}\label{thm:kSP-d}
Given any sequence $\requests$ (resp. $\requests'=\servers\requests$) of $n$ requests (resp. $n$ requests and $k$ initial server locations) in $d$ dimensions with a spread of $\Delta$, and a value $1 \le k \le n$, there exists a deterministic algorithm that computes the optimal solution for the instance of $k$-SP (resp. $k$-SPI) problem under the $\ell_p$ norm in $\tilde{O}(\min\{nk, n^{2-\frac{1}{2d+1}}\log \Delta\}\cdot\Phi(n))$ time.
\end{theorem}

% The next lemma then follows since each iteration of the extended Hungarian search procedure on $\cell$ takes $\tilde{O}(n_\cell\Phi(n))$ time.

% \begin{restatable}{lemma}{GlobalTime}\label{lemma:global-time}
%     For any cell $\cell$ of $\mcH$, the total execution time of the extended Hungarian search procedure on $\cell$ takes $\tilde{O}(n^{4/5}n_\cell\Phi(n))$ time.
% \end{restatable} 















